\section*{Chapitre \ref{ch:b3:learning-memory} --- Plasticité neuronale, apprentissage et mémoire}

\begin{solution}{exo:b3:learning-memory:1}
La \omterm{def:b3:learning-memory:kinds}{mémoire déclarative} (faits, événements) était abolie chez H.M. --- il
ne pouvait rien rappeler de nouveau --- tandis que la mémoire non
déclarative survivait~: il apprenait normalement le dessin en miroir
sans se souvenir des séances. Savoir-faire~: striatum et \omterm{def:b3:nervous-systems:plan}{cervelet}~; peur
conditionnée~: amygdale~; réflexes moteurs conditionnés~: \omterm{def:b3:nervous-systems:plan}{cervelet}~;
amorçage~: cortex sensoriel.
\end{solution}

\begin{solution}{exo:b3:learning-memory:2}
Quand une cellule présynaptique contribue à plusieurs reprises à faire
décharger une cellule postsynaptique, leur connexion se renforce. D'où
une PLT spécifique de l'entrée (seules changent les synapses actives,
qui y ont pris part), associative (une entrée faible active pendant que
la cellule décharge sous une forte est renforcée) et coopérative (assez
d'entrées doivent agir ensemble pour faire décharger la cellule).
\end{solution}

\begin{solution}{exo:b3:learning-memory:3}
Il ne conduit que lorsque le glutamate est lié (la cellule présynaptique
a déchargé) \emph{et} que la membrane est assez dépolarisée pour expulser
le magnésium (la cellule postsynaptique est active)~: les deux
conditions de la \omterm{def:b3:learning-memory:ltp}{règle de Hebb} dans une seule protéine. AP5 avant le
tétanos~: pas d'entrée de calcium, pas de PLT, bien que la transmission
par les \omterm{prop:b3:learning-memory:nmda}{récepteurs AMPA} soit normale. AP5 après~: une PLT déjà induite
n'est pas affectée --- le récepteur est nécessaire à l'induction, non au
maintien.
\end{solution}

\begin{solution}{exo:b3:learning-memory:4}
À court terme~: modification covalente de protéines existantes --- la
PKA phosphoryle et ferme un canal potassique, ce qui élargit le
potentiel et augmente la libération~; aucune protéine nouvelle. À long
terme~: la PKA active CREB, de nouveaux gènes sont transcrits et de
nouvelles terminaisons synaptiques poussent. Un inhibiteur de la
synthèse protéique donné pendant un entraînement répété laisse intacte
la facilitation à court terme et abolit le souvenir mesuré quelques
jours plus tard~; donné après, il ne fait rien.
\end{solution}

\begin{solution}{exo:b3:learning-memory:5}
$0.7^{n} \le 0.1$~: $n \ge 6.5$, donc $7$ essais. À partir de
$0.9\lambda$, $0.9\times 0.7^{n} < 0.1$~: encore $7$ essais. Symétriques
parce que le même $\alpha$ gouverne les deux sens. Chez l'animal
l'extinction est d'ordinaire plus lente et le souvenir initial revient
après un repos (le rétablissement spontané)~: l'extinction est un
apprentissage nouveau qui étouffe l'ancien, non son effacement.
\end{solution}

\begin{solution}{exo:b3:learning-memory:6}
Vecteurs propres $(1,1)/\sqrt{2}$ de valeur propre $1.5$ et
$(1,-1)/\sqrt{2}$ de $0.5$~; Hebb tourne $\mathbf{w}$ vers $(1,1)$, et le
neurone en vient à détecter les deux entrées déchargeant ensemble. Pas
d'Oja~: $y = 1$,
$\Delta\mathbf{w} = 0.1\times 1\times((1,1) - 1\times(1,0)) = (0,0.1)$,
donc $\mathbf{w} = (1, 0.1)$.
\end{solution}

\begin{solution}{exo:b3:learning-memory:7}
$500\times 0.05 = 25$ ions libres~; $25/6\times 10^{23} = 4.2\times 10^{-23}$
mol dans $10^{-16}$ L~: \qty{0.42}{\micro mol/L}, quatre fois le niveau
de repos mais au-dessous du $K_{d}$ de la calmoduline --- activation
partielle seulement~; il faut plusieurs ouvertures de canaux ensemble
pour atteindre la gamme micromolaire qui allume CaMKII.
\end{solution}

\begin{solution}{exo:b3:learning-memory:8}
A est entraîné jusqu'à $V_{A} = \lambda$~: la dopamine décharge en
bouffées aux premiers essais et se tait à mesure que la récompense
devient prédite. Lors des essais composés la prédiction
$V_{A} + V_{B} = \lambda$ est déjà exacte, l'erreur est nulle, la
dopamine se tait, et $V_{B}$ ne change jamais~: B est bloqué. Éprouvé
seul, B n'évoque aucune réponse~; si une récompense suit alors B seul,
elle est inattendue et la dopamine décharge en bouffée.
\end{solution}

\begin{solution}{exo:b3:learning-memory:9}
(a) Pas de PLT dans \omterm{def:b3:learning-memory:hippocampus}{CA1} et un déficit de mémoire spatiale, les autres
apprentissages étant intacts. (b) Mémoire à court terme normale, mémoire
à long terme absente à \qty{24}{h}. (c) Aucun effet --- la \omterm{def:b3:learning-memory:kinds}{consolidation}
est finie (sauf si le souvenir est réactivé, auquel cas il redevient
labile). (d) La souris s'immobilise dans le contexte sûr~: le souvenir
est rappelé artificiellement en réactivant les cellules qui l'encodaient.
\end{solution}

\begin{solution}{exo:b3:learning-memory:10}
$\Delta|\mathbf{w}|^{2} \approx 2\mathbf{w}\cdot\Delta\mathbf{w} = 2\eta
y^{2} \ge 0$~: la longueur ne décroît jamais et croît dès que le neurone
décharge. En un point fixe d'Oja $C\mathbf{w} = (\mathbf{w}^{\mathsf{T}}C
\mathbf{w})\mathbf{w}$, donc $\mathbf{w}$ est un vecteur propre avec $\lambda =
\lambda|\mathbf{w}|^{2}$, d'où $|\mathbf{w}| = 1$. Une croissance sans
borne saturerait toutes les synapses, abolirait la sélectivité et
mènerait à une excitation emballée~; les normalisateurs biologiques sont
la mise à l'échelle synaptique (toutes les synapses d'un neurone
réduites quand il est trop actif), la DLT, et la réduction des synapses
pendant le sommeil.
\end{solution}

\begin{solution}{exo:b3:learning-memory:11}
Rat~: $0.14\times 3\times 10^{5} \approx 4\times 10^{4}$~; humain~: $0.14
\times 2\times 10^{6} \approx 3\times 10^{5}$. L'estime suppose des
motifs aléatoires et denses et une connectivité complète~; les motifs
réels sont parcimonieux (ce qui élève beaucoup la capacité) et
structurés, et l'\omterm{def:b3:learning-memory:kinds}{hippocampe} est un magasin temporaire qui remet les
souvenirs à un cortex de $10^{10}$ neurones --- le magasin d'une vie est
le cortex, et le nombre hippocampique est une taille de tampon, non un
compte de souvenirs.
\end{solution}

\begin{solution}{exo:b3:learning-memory:12}
Les entrées de l'œil ouvert sont corrélées à la décharge de la cellule
corticale et sont renforcées~; celles de l'œil fermé ne le sont pas, et
sous la compétition pour un poids total limité elles s'affaiblissent ---
Hebb plus normalisation. Les deux yeux fermés, aucune entrée n'est
favorisée, de sorte que toutes deux gardent un territoire (faible) et
l'effet est plus doux. La période se ferme à mesure que les circuits
inhibiteurs mûrissent, que les filets périneuronaux se raidissent autour
des synapses, et que les sous-unités du \omterm{prop:b3:learning-memory:nmda}{récepteur NMDA} passent à une
forme moins permissive.
\end{solution}

\begin{solution}{pb:b3:learning-memory:1}
\textbf{1.} $10\times 10^{3} = 10^{4}$ ions~; \qty{2}{\%} libres, $200$~;
$200/6\times 10^{23} = 3.3\times 10^{-22}$ mol dans $8\times 10^{-17}$ L~:
\qty{4.2}{\micro mol/L}.
\textbf{2.} Oui, au-dessus du seuil de \qty{2}{\micro mol/L}. Deux
récepteurs~: $40$ ions libres, \qty{0.83}{\micro mol/L} --- en dessous.
\textbf{3.} À \qty{0.83}{\micro mol/L} la calcineurine est active et
CaMKII non~: les \omterm{prop:b3:learning-memory:nmda}{récepteurs AMPA} sont retirés et la synapse s'affaiblit
(DLT). Un seul ion, deux enzymes d'affinités différentes~: une montée
ample et rapide atteint la kinase, une montée faible et lente seulement
la phosphatase.
\textbf{4.} Une seconde entrée survenant en quelques constantes de temps
ajoute son calcium à ce qui reste de la première~; après \qty{50}{ms} il
ne reste presque rien. L'exigence de coïncidence du NMDA et la
décroissance en \qty{20}{ms} donnent ensemble une fenêtre de quelque
\qty{20}{ms} de part et d'autre --- la fenêtre temporelle des potentiels.
\textbf{5.} Les tampons capturent l'essentiel du calcium en une fraction
de micromètre et le col étroit de l'épine ralentit la diffusion vers la
dendrite, de sorte que la montée reste dans l'épine qui était active~;
une voisine à \qty{0.5}{\micro m} ne voit rien, et seule la synapse
active change.
\textbf{6.} $(70 - 30)/5 = 8$ entrées synchrones, ou un potentiel
d'action rétropropagé venu du corps cellulaire~: une synapse seule ne
peut débloquer ses propres \omterm{prop:b3:learning-memory:nmda}{récepteurs NMDA} --- la \omterm{def:b3:structural-biology:allostery}{coopérativité} est
inscrite dans la physique.
\textbf{7.} $(1,1)/\sqrt{2}$ de valeur propre $1.6$~; $(1,-1)/\sqrt{2}$
de $0.4$.
\textbf{8.} $C\mathbf{w}_{0} = (1, 0.6)$, $\mathbf{w}_{1} = (1.10,
0.06)$, angle à $(1,1)$~: $45^{\circ} - 3.1^{\circ} = 41.9^{\circ}$.
$\mathbf{w}_{2} = (1.21, 0.13)$~: $38.8^{\circ}$. $\mathbf{w}_{3} = (1.34,
0.22)$~: $35.8^{\circ}$.
\textbf{9.} $w_{2}/w_{1} \to 1$ tandis que $|\mathbf{w}|$ croît sans
borne (d'un facteur $1.16$ par pas).
\textbf{10.} $C\mathbf{w} = (1.28, 1.28)$~; $\mathbf{w}^{\mathsf{T}}C
\mathbf{w} = 2.05$~; $\Delta\mathbf{w} = 0.1\times(1.28 - 2.05\times 0.8)
= -0.036$ pour chaque composante~: $\mathbf{w} = (0.76, 0.76)$, qui se
réduit vers le vecteur unitaire $(0.71, 0.71)$.
\textbf{11.} Conjoint $y = 1.42$~; entrée 2 seule $y = 0.71$, la moitié.
\textbf{12.} L'entrée faible décharge pendant que la forte fait décharger
la cellule, elle est donc corrélée à la sortie et la \omterm{def:b3:learning-memory:ltp}{règle de Hebb} la
renforce --- le neurone apprend la cooccurrence.
\textbf{13.} $1 - 0.85^{n}$~: $0.56$, $0.80$, $0.96$.
\textbf{14.} $0.85^{n} \le 0.05$~: $n \ge 18.5$, donc $19$ essais.
\textbf{15.} Non~: la fraction de l'asymptote atteinte après $n$ essais
ne dépend pas de $\lambda$~; c'est l'asymptote qui double. Une plus
grande récompense peut élever $\alpha$ (la saillance), et un seuil
comportemental portant sur la valeur \emph{absolue} est atteint plus tôt.
\textbf{16.} $V_{B} = 0$~: bloqué. Le composé est entièrement prédit,
l'erreur est nulle, les neurones dopaminergiques se taisent.
\textbf{17.} B seul ne prédit rien, donc l'erreur vaut $\lambda$ et la
dopamine décharge en bouffée~; $V_{B}$ monte alors comme à
l'acquisition, $\lambda(1 -
0.85^{n})$.
\textbf{18.} La drogue délivre la bouffée qui signifie «~mieux que
prévu~» après quoi que ce soit qui l'ait précédée~; la règle incrémente
la valeur de ces indices et de ces actions à chaque occasion, si bien
qu'ils acquièrent une valeur indépendante de toute issue réelle --- la
quête devient compulsive.
\textbf{19.} $0.14\times 3\times 10^{5} \approx 4.2\times 10^{4}$.
\textbf{20.} $4.2\times 10^{4}/20 = 2100$ jours, environ six ans~; la
\omterm{def:b3:learning-memory:kinds}{consolidation} doit déplacer les souvenirs vers le cortex et l'oubli doit
libérer le magasin.
\textbf{21.} Écrire vite de nouveaux motifs dans le cortex écraserait
les synapses qui encodent les anciens (l'interférence catastrophique)~;
un \omterm{def:b3:learning-memory:hippocampus}{rejeu} lent et entrelacé laisse le cortex intégrer le nouveau à
l'ancien, tandis que l'\omterm{def:b3:learning-memory:kinds}{hippocampe} tient entre-temps le souvenir neuf ---
un tampon rapide qui alimente un magasin lent.
\textbf{22.} $0.03\times 3\times 10^{5} = 9000$ cellules par lieu. Des
motifs parcimonieux se recouvrent peu, de sorte qu'on peut en stocker
bien davantage avant que leur interférence ne corrompe la récupération~;
la capacité croît plus vite que $N$ à mesure que les motifs
s'appauvrissent, au prix d'une spécificité qui mobilise plus de cellules
par motif.
\textbf{23.} Compression d'un facteur $20$~: des cellules qui déchargent
à \qty{400}{ms} d'écart dans le parcours déchargent à \qty{20}{ms}
d'écart dans le \omterm{def:b3:learning-memory:hippocampus}{rejeu} --- à l'intérieur de la fenêtre temporelle. Tout
l'intérêt de la compression est d'amener à portée hebbienne des
séquences trop lentes pour être associées en temps réel.
\textbf{24.} À un jour le souvenir dépend encore de l'\omterm{def:b3:learning-memory:kinds}{hippocampe}~; à un
mois il a été consolidé dans le cortex et n'en a plus besoin ---
exactement la perte temporellement graduée de H.M., les souvenirs
récents disparus et les lointains épargnés.
\textbf{25.} Environ \qty{4.2}{\micro mol/L} de calcium libre après dix
ouvertures~; $19$ essais pour \qty{95}{\%}~; quelque $4\times 10^{4}$
motifs dans le \omterm{def:b3:learning-memory:hippocampus}{CA3} du rat.
\end{solution}
